%% ma: L11-B26-0004
%% lop: 11
%% chuong: 7
%% bai: 26
%% dang: 11.26.1
%% loai: TLN
%% muc: VD
%% dap_an: "15"
%% nguon: "(NBV)-ĐỀ SỐ 4-MỖI NGÀY 1 ĐỀ THI 2025.pdf (D04, MỖI NGÀY 1 ĐỀ - Đề số 4), Phần III Câu 3"
%% kien_thuc: "Bài 25 (hai mặt phẳng vuông góc, góc giữa hai mặt phẳng), Bài 26 (khoảng cách từ điểm đến mặt phẳng)"
%% trang_thai: chinh-thuc
%% ghi_chu: "Cùng dạng 11.26.1 (đổi điểm cần tính khoảng cách để quy về khoảng cách trong mặt đáy, nhân sin của góc giữa hai mặt phẳng).. Đáp án gốc 15 đúng (sympy: d = a căn 15 / 5)."
\begin{ex}
Cho hình chóp $S.ABCD$ có đáy $ABCD$ là hình thang vuông tại $A$ và $D$; $AB=AD=2a$; $DC=a$. Điểm $I$ là trung điểm của đoạn $AD$, hai mặt phẳng $(SIB)$ và $(SIC)$ cùng vuông góc với mặt phẳng $(ABCD)$. Mặt phẳng $(SBC)$ tạo với mặt phẳng $(ABCD)$ một góc $60^\circ$. Khoảng cách từ $D$ đến $(SBC)$ có dạng $\dfrac{a\sqrt n}{5}$, $n\in\mathbb N$. Tìm giá trị của $n$.
\shortans{15}
\loigiai{Hai mặt phẳng $(SIB),(SIC)$ cùng vuông góc với $(ABCD)$ và cắt nhau theo $SI$ nên $SI\perp(ABCD)$. Gọi $K$ là hình chiếu của $I$ trên $BC$ thì góc giữa $(SBC)$ và $(ABCD)$ là $\widehat{SKI}=60^\circ$.
Với mọi điểm $P$ thuộc $(ABCD)$ ta có $d(P,(SBC))=d(P,BC)\cdot\sin60^\circ$.
Đặt hệ trục trong $(ABCD)$: $A(0;0)$, $B(2a;0)$, $C(a;2a)$, $D(0;2a)$. Đường thẳng $BC$: $2x+y-4a=0$ nên $d(D,BC)=\dfrac{|2a-4a|}{\sqrt5}=\dfrac{2a}{\sqrt5}$.
Vậy $d(D,(SBC))=\dfrac{2a}{\sqrt5}\cdot\dfrac{\sqrt3}2=\dfrac{a\sqrt{15}}5$, suy ra $n=15$.}
\end{ex}
